% Curated from paper commit 7ee9157; original equations and protocol identities retained.
\section{Focal shunting and spatial transport}
\label{note:shunt_transport}

\subsection{Passive focal shunts attenuate a positive adjoint}

\begin{proposition}[No sign reversal in a passive grounded tree]
Let $G$ be a symmetric positive-definite nonsingular M-matrix, with
nonpositive off-diagonal entries. For a nonnegative adjoint source
$\bm s$, define $\bm q=G^{-1}\bm s$. Adding a
nonnegative shunt $\kappa$ at compartment $k$ gives
$G'=G+\kappa\bm e_k\bm e_k^{\mathsf T}$ and
\begin{equation}
\bm q'=\bm q-
\frac{\kappa q_k}{1+\kappa(G^{-1})_{kk}}G^{-1}\bm e_k.
\label{eq:s_sherman_shunt}
\end{equation}
Both $G^{-1}$ and $(G')^{-1}$ are elementwise nonnegative, hence
$0\leq\bm q'\leq\bm q$ componentwise. A focal passive shunt attenuates but
cannot reverse a positive adjoint field.
\end{proposition}

\noindent\emph{Proof.}
Equation~\ref{eq:s_sherman_shunt} is the Sherman--Morrison identity. A
nonsingular M-matrix has a nonnegative inverse; adding a nonnegative diagonal
term preserves the M-matrix property \citep{horn2013matrix}. The subtracted vector is nonnegative,
and $\bm q'=(G')^{-1}\bm s\geq0$. \hfill$\square$

This ordering does not by itself guarantee descendant selectivity relative to
a matched current injection; that spatial comparison depends on the Green's
function, focal site and electrotonic regime and is evaluated numerically in
the focal experiments.

\begin{corollary}[A focal shunt applies gains to its ancestry partition]
\label{cor:shunt_ancestry_gain}
Assume additionally that the nonzero off-diagonal entries of $G$ form a
connected tree rooted at the soma $s$, and that the adjoint source is
$\ell\bm e_s$ for a fixed scalar somatic loss derivative $\ell$. Write
$R=G^{-1}$ and $\bm h=R\bm e_s>0$. For a shunt at $k$, let
$\mathcal A_k$ be the vertices on the path from $s$ to $k$, including both
endpoints. For each $a\in\mathcal A_k$, define the disjoint block
$P_a=\{i:\operatorname{LCA}(i,k)=a\}$, where $\operatorname{LCA}$ is the
lowest common ancestor. Then
\begin{equation}
q'_i=\eta_a q_i\quad(i\in P_a),\qquad
\eta_a=1-\frac{\kappa R_{ks}R_{ak}}
{(1+\kappa R_{kk})R_{as}},\qquad
\eta_k=\frac{1}{1+\kappa R_{kk}}.
\label{eq:s_shunt_partition_gain}
\end{equation}
For finite $\kappa\geq0$, every gain satisfies $0<\eta_a\leq1$.
Let $B_k$ have columns $\bm 1_{P_a}$ and define the baseline-weighted
partition dictionary $D_k=\operatorname{diag}(\bm h)B_k$. With
$\Gamma_k=\operatorname{diag}(\eta_a)$,
\begin{equation}
\bm q=\ell D_k\bm 1,\qquad
\bm q'=\ell D_k\Gamma_k\bm 1.
\label{eq:s_shunt_partition_dictionary}
\end{equation}
Thus the adjoint effect is exactly a diagonal gain in this partition basis.
\end{corollary}

\noindent\emph{Proof.}
Removing a vertex $a$ separates a tree into components. If $a$ lies on the
path from $i$ to $j$, solving the source-free equations in the component
containing $i$, with boundary voltage at $a$, gives
$R_{ij}=R_{ia}R_{aj}/R_{aa}$. For
$a=\operatorname{LCA}(i,k)$, apply this identity to $(i,k)$ and $(i,s)$;
it also holds when an endpoint equals $a$. Hence
$R_{ik}/R_{is}=R_{ak}/R_{as}$. Substitution into
Eq.~\ref{eq:s_sherman_shunt} gives Eq.~\ref{eq:s_shunt_partition_gain}.
Positivity of both tree inverse matrices gives positive ratios, and the
preceding proposition bounds them by one. The formula also holds when
$\ell=0$, when both adjoints vanish. The disjoint blocks partition all
vertices, yielding Eq.~\ref{eq:s_shunt_partition_dictionary}.
\hfill$\square$

The block $P_k$ is the whole subtree below the shunt. For an ancestor $a$
above $k$, $P_a$ contains $a$ and its sister subtrees off the path to $k$.
If $c$ is the next vertex toward $k$, its indicator is
$\bm 1_{P_a}=\bm 1_{\mathrm{Desc}(a)}-\bm 1_{\mathrm{Desc}(c)}$,
whereas $\bm 1_{P_k}=\bm 1_{\mathrm{Desc}(k)}$. The gain field therefore
lies in the span of nested ancestry indicators. The diagonal form requires
the disjoint partition basis; an arbitrary selected set of overlapping
routes need not have diagonal gains. Even for the ancestry chain, a change
from partition to nested-subtree coordinates generally makes the gain
matrix non-diagonal, while preserving that full span.

For the synaptic gradient, the local driving force remains a separate factor:
$\gamma'_i=\eta_a q_i(E_{\mathrm E}-V'_i)$ for $i\in P_a$.
Its within-block variation is therefore not constrained to a single gain
relative to $\gamma_i=q_i(E_{\mathrm E}-V_i)$. The corollary describes
somatic-source adjoint transport on a passive tree, rather than the complete
synaptic update, a distributed-error source or a nonlinear finite-dose
response. It connects the morphology-defined dictionary to a conductance
gain without implying a learning advantage.

\begin{corollary}[Spectral size of a focal passive perturbation]
Under the preceding assumptions, let $\mathcal R=G^{-1}$ and
$\mathcal R'=(G+\kappa\bm e_k\bm e_k^{\mathsf T})^{-1}$. The inverse-operator change has
rank at most one, with rank one for $\kappa>0$, and
\begin{equation}
\lVert \mathcal R'-\mathcal R\rVert_2
=\frac{\kappa\lVert \mathcal R\bm e_k\rVert_2^2}
{1+\kappa \mathcal R_{kk}}
\leq
\frac{\kappa\lVert \mathcal R\rVert_2^2}
{1+\kappa/\lambda_{\max}(G)}.
\label{eq:s_focal_inverse_norm}
\end{equation}
Consequently, for any adjoint source $\bm s$,
$\lVert(\mathcal R'-\mathcal R)\bm s\rVert_2\leq
\lVert \mathcal R'-\mathcal R\rVert_2\lVert\bm s\rVert_2$.
\end{corollary}

\noindent\emph{Proof.}
Sherman--Morrison gives
$\mathcal R'-\mathcal R=-\kappa(\mathcal R\bm e_k)(\mathcal R\bm e_k)^{\mathsf T}/
(1+\kappa \mathcal R_{kk})$.
The norm of this symmetric outer-product matrix is its scalar coefficient times
$\lVert \mathcal R\bm e_k\rVert_2^2$. Use
$\lVert \mathcal R\bm e_k\rVert_2\leq\lVert \mathcal R\rVert_2$ and
$\mathcal R_{kk}\geq\lambda_{\min}(\mathcal R)=1/\lambda_{\max}(G)$ for the bound.
\hfill$\square$

Equation~\ref{eq:s_focal_inverse_norm} provides only a global magnitude ceiling
on the total adjoint change. It does not predict the localization contrast or
the observed crossover with membrane conductance. Those spatial effects depend
on the Green's-function profile and electrotonic length constant of the cable.

\subsection{Focal shunting perturbations}
\label{note:focal}

Route capacity does not show that conductance can selectively alter transport.
We therefore compared a focal shunt with a baseline-current-matched injection
at the same site of each reciprocal reconstructed cable and quantified the
resulting change in the exact modeled gradient.

\subsection{Adjoint formulation and baseline-current-matched current injection}

For passive conductance matrix $G$ and current-source vector $\bm b$, the
voltage vector was $\bm V=G^{-1}\bm b$. Let $\bm e_s$ select the somatic
coordinate and let $w_i$ denote the excitatory conductance at site $i$. For a
loss depending on somatic voltage, the adjoint error and conductance gradient
were
\begin{equation}
\bm q=G^{-1}\bm e_s\frac{\partial\mathcal L}{\partial V_s},
\qquad
\frac{\partial\mathcal L}{\partial w_i}=q_i(E_{\mathrm E}-V_i).
\end{equation}
Adding focal inhibitory conductance \(\kappa_k\) at node $k$, selected by
the coordinate vector $\bm e_k$, gives
\[
G'=G+\kappa_k\bm e_k\bm e_k^{\mathsf T},
\qquad
\bm b'=\bm b+\kappa_kE_{\mathrm I}\bm e_k .
\]
The shunt therefore changes both conductance loading and the forward source.
Holding the somatic loss derivative fixed, the Sherman--Morrison identity
gives the transport change
\begin{equation}
\bm q'=\bm q-
\frac{\kappa_kq_k}{1+\kappa_k(G^{-1})_{kk}}G^{-1}\bm e_k.
\label{eq:s_sherman}
\end{equation}
The reversal \(E_{\mathrm I}\) does not appear in
Eq.~\ref{eq:s_sherman} because the adjoint transport depends on \(G'\), not on
the forward source \(\bm b'\). It remains in the solved voltage and in every
inhibitory driving force.
Corollary~\ref{cor:shunt_ancestry_gain} supplies the exact partition-gain form. The full synaptic gradient additionally changes through its local driving force.

Let $r_k=(G^{-1})_{kk}>0$. Then
\begin{equation}
\Delta\bm q=-a_kG^{-1}\bm e_k,
\qquad
a_k=\frac{\kappa_kq_k}{1+\kappa_kr_k},
\qquad
\frac{\partial |a_k|}{\partial\kappa_k}
=\frac{|q_k|}{(1+\kappa_kr_k)^2}\geq0.
\label{eq:s_shunt_amplitude}
\end{equation}
The derivative is strictly positive when $q_k\neq0$; if $q_k=0$, the perturbation vanishes. Its amplitude saturates at $|q_k|/r_k$. For a nonzero perturbation, the normalized spatial profile is fixed by the Green's-function column $G^{-1}\bm e_k$ in this linear passive state. For a
descendant set $D(k)$ and a prespecified matched comparison set $C(k)$, define
\begin{equation}
S_k(D,C)=
\frac{\operatorname{median}_{i\in D(k)}|(G^{-1})_{ik}|}
{\operatorname{median}_{j\in C(k)}|(G^{-1})_{jk}|}.
\label{eq:s_electrotonic_selectivity}
\end{equation}
For $\kappa_k>0$, $q_k\neq0$ and a positive comparison median, the common amplitude $|a_k|$ cancels. Under these conditions, $S_k>1$ is necessary and sufficient for a positive descendant-minus-comparison contrast in absolute adjoint change at fixed $k$. It is only a transport-level criterion. Relative
or logarithmic changes additionally divide by the baseline $q_i$, and the full
synaptic gradient also contains the post-perturbation driving force. This is
why axial coupling, leak, background conductance and local input resistance
must be varied explicitly rather than inferred from topology alone.
The baseline-current-matched current injection applied
$\kappa_k(E_{\mathrm I}-V_k)$ at the same node without changing $G$. This equals the
current carried by the added shunt at the unperturbed voltage; its local
voltage effect matches the shunt only to first order in dose. An independently
solved state-matching compensatory current injection restored baseline somatic
voltage after both interventions. This is not a continuously adjusted voltage
clamp. The somatic scalar error was therefore fixed. Shunting changed adjoint
transport, whereas the current-injection control left it unchanged. Dendritic voltages
were not clamped, so both interventions could also change local driving forces.

\subsection{Prespecified endpoint and numerical validation}

Before running the analysis, we fixed the site-selection rule, primary dose,
comparison groups, and localization endpoint. Focal sites were selected by a
depth-stratified rule without access to gradient outcomes. A site required at
least three excitatory descendants and
three comparison sites. At most 16 eligible sites were sampled per cell across
quartiles of topological depth. The primary model used excitatory and
inhibitory synaptic scales of 0.35 and normalized reversal potentials
$E_{\mathrm E}=1$ and $E_{\mathrm I}=-0.2$. Synapses were classified as descendants, ancestors,
sister-subtree sites or unrelated sites. Unrelated synapses within each
template were matched to descendant cable depth. For baseline gradient $g_i$
and perturbed gradient $g_i'$, we defined
$m_i=\left|\log[(|g_i'|+\varepsilon)/(|g_i|+\varepsilon)]\right|$, where
$\varepsilon=10^{-15}\max(1,\max_i|g_i|)$. For each site, the localization
index was the median $m_i$ among descendants minus the median among
depth-matched unrelated synapses. Site effects were first averaged within cell.

The somatic loss was one-half the squared difference between somatic voltage
and a target one normalized unit below its baseline value. Conductance doses
were 0.25, 0.5, 1, and 2 times the local leak-plus-synaptic conductance. The
state-matching current was solved after each perturbation and then treated as
fixed when evaluating the post-intervention conductance gradient. The topology
control sampled 200 relation templates from other focal sites, stratified by
focal-site topological-depth quartile when another site was available;
otherwise it sampled another site in the same cell.

At unit dose, 101 sites in eight cells gave mean localization 0.104 for shunting and 0.035 for current injection. Their paired cell-level contrast was 0.069 log units (95\% interval, 0.053--0.086), positive in all eight cells. True relations also exceeded the depth-stratified relation control. The contrast grew with relative dose; Source Data retain every site, dose and template.

We separated the two factors in the exact gradient algebraically. Writing
$d_i=E_{\mathrm E}-V_i$ and $d'_i=E_{\mathrm E}-V'_i$, full shunting uses $q'_id'_i$,
adjoint-only substitution uses $q'_id_i$, and driving-force-only substitution
uses $q_id'_i$. At unit dose, their mean localization indices were 0.104,
0.130 and 0.026, respectively, compared with 0.035 for the current-injection
control. Adjoint-only minus driving-force-only localization was 0.104 (95\%
interval 0.080--0.125), positive in all eight cells.

Full shunting and the driving-force-only state provide a direct comparison
of adjoint transport: they share the same post-shunt voltage field and differ
only in whether the adjoint field is updated. Their localization difference was 0.079 and was positive in all
eight cells. This is the exact adjoint-replacement contrast plotted in main
Fig.~8C.

The gradient factors and localization statistic are nonlinear, so these four
outcomes do not form an additive decomposition. To allocate the
full-shunt-minus-current-injection contrast, we used two-factor Shapley values.
For $v_{00}$ equal to current-injection localization, $v_{10}$ adjoint-only,
$v_{01}$ driving-force-only and $v_{11}$ full-shunt localization,
\begin{align}
\phi_q&=\tfrac12[(v_{10}-v_{00})+(v_{11}-v_{01})],\\
\phi_d&=\tfrac12[(v_{01}-v_{00})+(v_{11}-v_{10})].
\end{align}
The allocations sum exactly to the full-shunt-minus-current contrast. Mean adjoint and driving-force allocations were 0.087 and $-0.017$, respectively: driving-force change partly cancelled the adjoint contribution. This is an algebraic allocation of a nonlinear statistic; the current-matched control is not a literal no-factor corner of a physical factorial. Repeating the allocation with the unperturbed field as its reference is retained in Source Data. Factor substitutions are not independent biological interventions.

All conductance matrices in the primary analysis were positive definite. The
somatic state-matching residual was numerically zero. Analytic adjoint gradients agreed
with finite differences to a maximum relative error of
$1.03\times10^{-6}$. The Source Data indexed in Table~\ref{tab:source_index} retain prespecified
conductance-scale and reversal-potential sensitivity runs. The direction of
the shunting-minus-current-injection effect was positive in every condition,
with seven or eight positive cells. These results establish the mechanism in a passive model on measured trees;
they do not establish endogenous inhibitory teaching signals in vivo.

With proxy-classified contacts removed, the primary
shunt-minus-current-injection contrast was 0.101 (0.075--0.129) across 55 sites
and was positive in all eight cells. In the independent, non-overlapping v661
sample, 45 cells supplied 235 eligible sites. The same contrast was 0.138
(0.116--0.162), positive in 45/45 cells. Forty cells
with at least two eligible sites supplied the reassigned-relation control; the
true relation exceeded this control by 0.118 (0.094--0.144), with
39 positive cells and one numerical tie (Supplementary
Fig.~\ref{fig:si_shunt_replication}A,B). Cells, rather than
sites or relation draws, were the
units of inference.

To test sensitivity to electrical calibration, we also computed axial and
leak conductances in physical units from the same compressed geometry. For
segment radius $r$, length $\ell$, axial resistivity $R_a$ and specific
membrane resistance $R_m$, a cylindrical non-root segment had axial conductance
$\pi r^2/(R_a\ell)$ and membrane leak $2\pi r\ell/R_m$; the root was assigned
spherical area $4\pi r^2$. All conductances were converted to nanosiemens (nS). Because MICrONS
synapse size is not an electrical conductance, synaptic sizes retained the
prespecified relative scaling with the median physical leak as reference. We
varied $R_a$ and $R_m$ while retaining the site-selection rule, relative dose, gradient definition and current-injection control. For each condition and intervention, we re-solved the compensating somatic current to restore that condition's baseline somatic voltage. The Source Data indexed in Table~\ref{tab:source_index} retain this electrotonic dependence. The large contrast observed under the normalized proxy was recovered only in
low-$R_m$, high-conductance settings; at
$R_a=150\,\Omega\,\mathrm{cm}$ and
$R_m=15{,}000\,\Omega\,\mathrm{cm}^2$ it was negligible or small. These are
sensitivity settings, not fits to the recorded cells, and the model remains
passive and steady-state.

A broader parameter matrix, specified before numerical validation, tested
membrane resistance, background conductance and dose definition jointly. It used
\(R_a=150\,\Omega\,\mathrm{cm}\) and crossed
\(R_m\in\{300,1000,15{,}000\}\,\Omega\,\mathrm{cm}^2\), distributed leak-like
background conductance of 0, 1 or 4 times baseline leak, and three dose
definitions. Relative-local doses were 0.25, 1 and 2 times membrane-plus-
synaptic conductance; fixed doses were 0.05, 0.5 and 5 nS; input-normalized
doses were 0.25, 1 and 4 times $[(G^{-1})_{kk}]^{-1}$. The same 101 sites in
eight cells supplied 16,362 site--condition--intervention rows. Every
conductance matrix was positive definite (minimum eigenvalue 0.107 nS) and the
maximum somatic-state residual was $3.33\times10^{-16}$.

At fixed absolute dose, the conductance-specific contrast did not follow the
raw focal-shunt localization alone. At 0.5 nS, shunt-minus-current-injection
localization was 0.0104 (0.0080--0.0126), 0.0073 (0.0044--0.0104), and
$-0.0039$ ($-0.0075$ to $-0.0004$) at $R_m=300$, 1,000 and 15,000,
respectively. At 5 nS the corresponding values were 0.0974, 0.0741 and 0.0061;
the last interval included zero. Distributed background conductance restored
the input-normalized effect in the high-resistance condition: at unit
input-conductance dose, the contrast changed from $-0.0004$ with no background
to 0.1210 and 0.2587 at background multipliers one and four. Across all 8,181
focal-shunt rows, every descendant gradient magnitude was attenuated, none was
enhanced and no sign flip occurred. The passive mechanism therefore selectively attenuated gradients; sign
reversal would require dynamics or active conductances absent from this phase (Supplementary Fig.~\ref{fig:si_shunt_sensitivity}).

\subsection{Weak-channel steady-state linearization check}

We checked the local-Jacobian calculation after adding weak
voltage-dependent currents to the high-conductance cable.
The physical reconstructed-tree system used
$R_a=150\,\Omega\,\mathrm{cm}$,
$R_m=1{,}000\,\Omega\,\mathrm{cm}^2$ and distributed background conductance
equal to baseline leak. Steady-state sodium (Na), potassium (K), calcium (Ca),
hyperpolarization-activated cyclic nucleotide-gated (HCN) and
N-methyl-D-aspartate receptor (NMDA) current families had fixed reversal
potentials, half-activation voltages, slopes and gate
powers specified in the predefined parameter file. Global and segment-level
maximal conductances were drawn independently from prespecified log-normal
distributions. Median maximal Na, K, Ca and HCN conductances were 0.06,
0.10, 0.025 and 0.05 times baseline leak, respectively; NMDA conductance was
0.15 times the excitatory conductance. Gate activation further scales these
currents. These prespecified weak-channel draws are not cell-specific
density estimates or a test of strongly regenerative operation.

For each draw, a residual-reducing Newton solve found the nonlinear steady
state. The predefined acceptance criteria required all voltages to lie in the
specified range, maximum absolute state residual below $10^{-10}$, and minimum
eigenvalue of the symmetrized local Jacobian above $10^{-8}$. The implicit-adjoint calculation
then used the exact inverse of this baseline Jacobian. For channel family $a$, let $\bar g_a$ be maximal conductance, $h_a(V)$ the
steady-state gating factor and $E_a$ the reversal potential. The current
$I_a(V)=\bar g_a h_a(V)(V-E_a)$ contributed the diagonal term
$\bar g_a[h_a(V)+h_a'(V)(V-E_a)]$; thus the baseline Jacobian includes
voltage-dependent gate derivatives. Focal shunting and a
baseline-current-matched current injection were evaluated at 0.25, 1 and 4
times local input conductance within this fixed local linearization.
Sherman--Morrison updated the Jacobian inverse for the added focal
conductance, and a compensating somatic current restored baseline somatic
voltage in that linearized system. The active gates and their derivatives
were not re-evaluated at a new nonlinear equilibrium after each finite dose.
Up to 256 attempts were permitted to obtain 64
accepted draws per cell. All first 64 attempts were accepted in every cell,
giving 512 equilibria, 101 focal sites and 38,784 site--draw--condition rows.
All values were finite, the largest equilibrium residual was
$9.51\times10^{-11}$ and the smallest Jacobian eigenvalue was 1.084.

Channel draws and focal sites were averaged within each of the eight cells
before inference. At doses 0.25, 1 and 4, mean shunt localization was 0.159,
0.529 and 1.339, respectively (Table~\ref{tab:source_index}). At unit dose,
shunting exceeded the current-injection control by 0.382 (95\% cell-bootstrap
interval 0.356--0.409), positive in 8/8 cells (two-sided Wilcoxon $P=0.0078$).
Every modeled descendant gradient was attenuated, none was enhanced or changed
sign, and descendant gradient energy fell to 0.330 of baseline. The dose
response closely followed the passive response at the same electrical
calibration. This agreement checks the weak-channel linearization; it does
not establish robustness to strongly active dendrites, channel kinetics or
spiking.

\subsection{Driving-force cancellation and the ancestry dependence of passive fields}
\label{sec:passive_field_diagnostics}
We reanalyzed the eight original cable operators used for the common-broadcast comparison. For each focal site, the perturbation was the original dose 0.001 with somatic voltage restored. The finite log-gradient difference divided by dose separates into an adjoint term and a driving-force term. Energy uses the original excitatory-area site weights and sums over focal perturbations. The full focal ancestry partition is defined separately for each shunt and is distinct from the seven selected routes in the budgeted anatomical dictionary.

Driving-force energy was 0.990--3.496 times full-field energy across cells; these are ratios of squared magnitudes, not amplitudes. Weighted centered correlations with the adjoint term were $-0.979$ to $-0.913$. The negative cross term explains why large components yield a smaller full response. The full response retained 0.9711--0.9993 of its energy in the focal partition. Reconstruction agrees with the original operators to $10^{-8}$; numerical decomposition closes to rounding error. Cellwise energies, correlations and partition captures are supplied in Source Data.

The passive corollary depends on a tree Green's function, a somatic adjoint source and a diagonal conductance change. Choosing other passive perturbations under the same assumptions does not make the anatomical field ancestry-neutral. The budgeted dictionary comparison measures compression within this response family. In the disjoint 47-cell cohort, actual ancestry exceeds the mean degree--depth surrogate in 38 cells, but in eleven cells more than half of 200 surrogates match or exceed the actual tree. This descriptive count uses tolerance $10^{-12}$ and retains the full within-cell distribution. A different error source, active dynamics or measured fields would be required to test beyond the present structural dependence.
